	\chapter{Linear Algebra Formalism of Quantum Mechanics}
	
	\section{Vectors in Hilbert Space}
	
	\emph{Hilbert space} is defined as a complete inner product space 
	where all the square-integrable functions are represented by \emph{kets} $|f\rangle$.\\
	%
	The inner product of two kets, representing the product of projection lengths, is denoted as $\langle f | g \rangle$.\\
	%
	In analysis, the \emph{bra} $\langle f |$ stands for an integration operation:
	%
	\[
		\langle f | = \int\! f^{*}(x)\cdots\text{d}x
	\]
	Which then gives the analytical expression for inner products:
	\[
		\langle f | g \rangle = \int\! f^{*}(x)g(x)\,\text{d}x
	\]
	In linear algebra, 
	vectors are expanded in their coordinate form with respect to a certain basis.
	\[
		| f \rangle = 
		\begin{pmatrix}
			f_1 \\ f_2 \\ \vdots \\ f_n
		\end{pmatrix},
		\quad 
		|g\rangle =
		\begin{pmatrix}
			g_1 \\ g_2 \\ \vdots \\ g_n
		\end{pmatrix}
	\]
	%
	The bra notation $\langle f |$ is the transpose conjugate of ket, i.e.
	\[\langle f | = | f \rangle^{T*} = (f_{1}^{*}\ f_{2}^{*}\ \cdots \ f_{n}^{*}).\]
	Hence $\langle f | g \rangle = f_{1}^{*}g_1 + f_{2}^{*}g_2 + \cdots + f_{n}^{*}g_n$.\\ \\
	Two important properties of inner products:\\
	(1) Conjugate commutation: $\langle f | g \rangle = \langle g | f \rangle^{*} \qquad$ \\       
	(2) Schwarz inequality: $|\langle f | g \rangle|^{2} \leq \langle f | f \rangle \langle g | g \rangle$\\ \\
	For an arbitrary group of functions $\{|f_n\rangle\}$: \\
	If $\langle f_m | f_n \rangle = \delta_{mn}$ (discrete), 
	or $\langle f_n | f_{n'} \rangle = \delta(n-n')$ (continuous), 
	we say $\{|f_n\rangle\}$ is \emph{orthonormal}. 
	If $\{|f_n\rangle\}$ spans the whole Hilbert space, we say it's \emph{complete}.\\ \\
	Therefore, for an orthonormal basis $\{|e_n\rangle\}$, any vector $|\alpha\rangle$ can be expressed in its coordinate form:
	\[|\alpha\rangle = \sum_{n}c_n|e_n\rangle\ \ \text{(discrete) or }|\alpha\rangle = \int\! c_n|e_n\rangle\,\text{d}n \ \ \text{(continuous)},\]
	where the coefficients $c_n$ is the projection length of $|\alpha\rangle$ on each $|e_n\rangle$, i.e.
	\[c_n = \langle e_n | \alpha \rangle.\]
	\emph{Note that this is also a good way to understand Fourier transform, as the trig functions can be orthonormalized.}\\ \\
	In quantum mechanics, the \emph{quantum state} is described as a coordinate-free state vector $|S\rangle$ in Hilbert space, 
	and the \emph{wavefunction} is its component on the position eigenbasis:
	\[\psi(x) = \langle x | S \rangle,\]\label{1.1}
	where $|x\rangle$ is the eigenvector of position operator $\hat{x}$ with an eigenvalue of $x$.\\
	Similarly, there exist the momentum-space wavefunctions and the momentum eigenfunctions under position eigenbasis:
	\[\phi(p) = \langle p | S \rangle, \quad f_p(x) = \langle x | p \rangle,\]
	where $|p\rangle$ is the eigenvector of momentum operator $\hat{p}$ with eigenvalue $p$.\\
	
	\section{Operators}
	\subsection{Operators and Matrices}
	\subsubsection{Matrix Elements}\label{matrix element}
	
	\emph{Operators} are linear transformations, i.e. matrices, in Hilbert space, which take a vector as input and output another:
	\[|\beta\rangle = \hat{Q}|\alpha\rangle.\]
	Operators can be written in their coordinate form under an orthonormal basis like matrices with their matrix elements:
	\[|\alpha\rangle = \sum_{n}a_n|e_n\rangle, \ \ |\beta\rangle = \sum_{n}b_n|e_n\rangle.\]
	Similar to matrix-vector multiplication:
	\[\begin{pmatrix}
		\cdots\\\cdots\\b_m\\\cdots
	\end{pmatrix} = 
	\begin{pmatrix}
		\cdots & \cdots & \cdots & \cdots \\ 
		\cdots & \cdots & \cdots & \cdots \\
		Q_{m1} & Q_{m2} & \cdots & Q_{mn} \\
		\cdots & \cdots & \cdots & \cdots
	\end{pmatrix}
	\begin{pmatrix}
		a_1\\a_2\\\cdots\\a_n
	\end{pmatrix} \Rightarrow
	b_m = \sum_{n}Q_{mn}a_n.\]
	On the other hand, \[ \sum_{n} b_n|e_n\rangle = \sum_{n} a_n\hat{Q}|e_n\rangle.\]
	Multiply by $\langle e_m | $ on both sides and using orthonormality $\langle e_m | e_n \rangle = \delta_{mn}$:
	\[\sum_{n} a_n\langle e_m | \hat{Q} | e_n \rangle = \sum_{n} b_n \langle e_m | e_n \rangle = b_m\]
	Which gives the expression for \emph{matrix elements}:
	\[Q_{mn} = \langle e_m | \hat{Q} | e_n \rangle\]
	\subsubsection{Commutators}
	Just like matrices, operators may not commute, so we define the \emph{commutator} of two operations:
	\[[ \hat{A}, \hat{B}] = \hat{A}\hat{B} - \hat{B}\hat{A}\]
	If two operators commute, i.e. their commutator is zero, they are guaranteed to share a set of complete eigenbasis of Hilbert space.\label{commutator}
	\subsubsection{Eigenvectors and Eigenvalues}
	If the result of an operator on a vector is only scaling, 
	the vector is called the \emph{eigenvector} (or eigenstate, eigenfunction) of the operator, 
	and the scaling factor is called the \emph{eigenvalue} of this eigenvector, i.e. if:
	\[\hat{Q}\psi = q\psi,\]
	$\psi$ is the eigenvector of $\hat{Q}$ with its eigenvalue $q$.\\
	For example, we want to find the eigenfunction of the momentum operator:
	\[\hat{p}f_p(x) = -i\hbar\frac{\partial}{\partial x}f_p(x) = p\cdot f_p(x)\]
	One can easily notice that $f_p(x)$ is some exponential function $Ae^{ipx/\hbar}$.\\
	Although $f_p(x)$ is not square-integrable, we can still orthonormalize it via the criterion in continuous cases:
	\[\langle f_p|f_p' \rangle = |A|^2\int\!e^{i(p-p')x/\hbar} \,\text{d}x = |A|^2 \cdot 2\pi\hbar\delta(p-p').\]
	Hence we choose $A=1/\sqrt{2\pi\hbar}$ and the momentum eigenfunction with eigenvalue $p$ under position eigenbasis is:
	\[f_p(x) = \frac{1}{\sqrt{2\pi\hbar}}e^{ipx/\hbar}\]
	
	\subsection{Projection Operators and Identity Operators}\label{identity}
	For a normalized (unit) vector $|\alpha\rangle$, 
	the inner product $\langle \alpha | f \rangle $ is the projection length of $|f\rangle$ on $|\alpha\rangle$, 
	while the \emph{projection operator} is a mapping to the one-dimensional space spanned by $|\alpha\rangle$ which outputs a vector:
	\[\hat{P} = |\alpha\rangle\langle\alpha|\]
	Therefore, for a complete orthonormal basis $\{|e_n\rangle\}$, 
	the sum (or integral in the continuous case) is the \emph{identity operator} 
	which first separate out all the components and then add them back together, which implies the introduction of coordinates.
	\[\hat{I} = \sum_{n}|e_n\rangle\langle e_n| \ \ \text{or  }\hat{I} =\! \int|e_n\rangle\langle e_n|\,\text{d}n\]
	
	\subsection{Hermitian Operators}
	\subsubsection{Adjoint Operators}
	If two operators $\hat{A}$ and $\hat{B}$ satisfies
	$\langle \hat{A}f|g\rangle = \langle f|\hat{B}g\rangle$, \\
	$\hat{A}$ is then defined as the \emph{adjoint operator} of $\hat{B}$, denoted as $\hat{A} = \hat{B}^{\dag}$\\ \\
	Because $\langle \hat{Q}^{\dag}f | g \rangle = \langle f|\hat{Q}g\rangle$ 
	and $\langle \hat{Q}^{\dag}f | g \rangle = \langle g| \hat{Q}^{\dag}f \rangle^{*}$\\
	Therefore $\langle g| \hat{Q}^{\dag}|f \rangle = \langle f|\hat{Q} |g\rangle^{*}$, i.e. the matrix elements $Q^{\dag}_{gf} = Q^{*}_{fg}$.\\
	Therefore the adjoint operator is the \emph{transpose conjugate} of the original one:
	\[\hat{Q}^{\dag} = \hat{Q}^{T*}\]
	
	\subsubsection{Hermitian Operators}\label{calc}
	Operators such that $\hat{H} = \hat{H}^{\dag}$ are called \emph{Hermitian operators}. Hermitian operators possess the following properties:
	\[\langle g| \hat{H} |f \rangle = \langle f|\hat{H} |g\rangle^{*} \Leftrightarrow \langle \hat{H}f|g\rangle = \langle f|\hat{H}g\rangle\]
	Proof: Hermitian operators have real eigenvalues and orthogonal eigenvectors:\\
	Assume $|f\rangle$ and $|g\rangle$ are two distinct eigenvectors of the Hermitian operator $\hat{H}$ with eigenvalues $f\neq g$, \\
	\[\langle f| \hat{H} |f \rangle = \langle f|\hat{H} |f\rangle^{*} \Rightarrow f\langle f|f \rangle = f^{*}\langle f|f\rangle, \ f = f^*\]
	 hence eigenvalue $f$ is real.\\
	\[\langle g| \hat{H} |f \rangle = \langle f|\hat{H} |g\rangle^{*} 
	\Rightarrow g\langle f|g\rangle = f^*\langle g|f\rangle^* = f\langle f|g\rangle\]
	Therefore $\langle f|g\rangle=0$ hence eigenvectors are orthogonal. \\ \\
	Based on Hermitian operators there lies an \emph{anxiom} of quantum mechanics:\\
	All \emph{observables} are repersented by corresponding Hermitian operators constructed from the two postulates below, 
	and their eigenvectors can form a set of complete orthonormal eigenbasis.\\
	The two \emph{postulates} are of the position operator and momentum operator under position eigenbasis:
	\[\hat{x} = x, \ \ \hat{p_x} = -i\hbar\frac{\partial}{\partial x}.\]
	
	\subsection{Expectation Values}\label{expectation}
	For the eigenbasis $\{|a_n\rangle\}$ of the Hermitian operator $\hat{A}$ and their eigenvalues $a_n$, 
	the identity operator is $\hat{I} = \sum|a_n\rangle\langle a_n|$. Therefore:
	\[\hat{A}|\psi\rangle = \hat{A}\hat{I}|\psi\rangle = \sum_{n}\hat{A}|a_n\rangle\langle a_n|\psi\rangle = \sum_{n}a_n|a_n\rangle\langle a_n|\psi\rangle.\]
	Thus we obtain the \emph{spectral decomposition} of the operator $\hat{A}$:
	\[\hat{A} = \sum_{n}a_n|a_n\rangle\langle a_n|.\]
	When measuring the observable $A$ of the normalized state $|\psi\rangle$:
	\[|\psi\rangle = \sum_{n}c_n|a_n\rangle \ \ \text{   where   }\ \  \sum_{n}|c_n|^2 = 1.\]
	The probability of $|\psi\rangle$ \emph{collapsing} into one of the eigenstates $|a_n\rangle$ and thus getting the result of eigenvalue $a_n$ is $|c_n|^2$, 
	hence the \emph{expectation value} is:
	\[\langle A \rangle = \sum_{n}|c_n|^2a_n.\]
	Utilizing the spectral decomposition, we can calculate the expectation value:
	\[\langle\psi|\hat{A}|\psi\rangle = \langle A \rangle = \sum_{n}|c_n|^2a_n.\]
	
	\section{Change of Basis}
	As described in \ref{identity}, identity operators are based on different eigenbases, 
	so we can transform the abstract state $|S\rangle$ into coordinate form, for example:
	\[|S(x)\rangle = \hat{I}_{x}|S\rangle = \int\!\ket{x}\expect{x|S}\,\text{d}x = \int\!\psi(x)|x\rangle\,\text{d}x\]
	Accoding to this we can change functions between different bases using the abstract state $|S\rangle$ as a bridge. 
	For instance, we want to find the relation between position-space wavefunction and momentum-space wavefunction:
	\[\begin{aligned}
		\phi(p) & = \langle p|S\rangle = \langle p|\hat{I}_x|S\rangle \\
		& = \langle p | \int\!\psi(x)|x\rangle\,\text{d}x 
		 = \int \!\langle p|x \rangle \psi(x)\,\text{d}x \\
		& = \int \!\langle x|p \rangle^* \psi(x)\,\text{d}x 
		 = \int \!f_{p}^{*}(x)\psi(x)\,\text{d}x .
	\end{aligned}\]
	Similarly, we can change operators between different bases by projecting the operation result onto a different basis.
	For instance, we want to express the position operator $\hat{x}$ in the momentum eigenbasis $\{|p\rangle\}$:
	\[\begin{aligned}
		\langle p|\hat{x}|S \rangle
		& = \int\!\langle p|\hat{x}|x\rangle\psi(x)\,\text{d}x \\
		& = \int\!x\langle p|x\rangle\psi(x)\,\text{d}x \\
		& = \int\!x\frac{1}{\sqrt{2\pi\hbar}}e^{-ipx/\hbar}\,\text{d}x.
	\end{aligned}\]
	Using the property of exponential functions, we can take $x$ out of the integral:
	\[\langle p|\hat{x}|S \rangle = i\hbar\frac{\partial}{\partial p}\int\!\langle p|x\rangle\psi(x)\,\text{d}x = i\hbar\frac{\partial}{\partial p}\phi(p).\]
	Therefore under the momentum eigenbasis, the position operator is:
	\[\hat{x} = i\hbar\frac{\partial}{\partial p}.\]
	
	\section{Uncertainty Principles}
	\subsection{Heisenberg Uncertainty Principle}
	We can use the \emph{variance} $\sigma_{Q}^{2}$ to measure the uncertainty of an observable $Q$:
	\[\sigma_{Q}^{2} = \left\langle\left(\hat{Q} - q\right)^{2}\right\rangle = \left\langle\psi\left|\left(\hat{Q}-q\right)^{2}\right|\psi\right\rangle,\]
	where $q=\langle Q \rangle$ is some real number, so the operator $\left(\hat{Q}-q\right)$ is also Hermitian. Therefore:
	\[\sigma_{Q}^{2} = \left\langle\left(\hat{Q} - q\right)\psi\left|\left(\hat{Q} - q\right)\psi\right\rangle\right..\]
	If $\sigma_{Q}^{2}=0$, then $\left(\hat{Q} - q\right)\psi=0$ which means that the \emph{determined state} is the \emph{eigenstate} of an observable. 
	But if not, we consider two observables $A$ and $B$. Denote:
	\[f = \left(\hat{A} - a\right)\psi\ \ \text{    and    }\ \ g = \left(\hat{B} - b\right)\psi.\]
	Invoking the Schwarz inequality and the property of complex numbers that
	\[|z|^{2} = \text{Re}^2(z)+\text{Im}^2(z) \geq\text{Im}^2(z)= \left(\frac{1}{2i}(z-z^*)\right)^2,\]
	we can obtain the inequality that
	\[\sigma^{2}_{A}\sigma^{2}_{B} = \langle f|f \rangle\langle g|g \rangle	 
	\geq \left|\langle f|g \rangle\right|^2 \geq \left(\frac{1}{2i}\left(\langle f|g \rangle-\langle g|f \rangle\right)\right)^2.\]
	Evaluate $\langle f|g \rangle$:
	\[\begin{aligned}
		\langle f|g \rangle
		& = \left\langle\left(\hat{A} - a\right)\psi\left|\left(\hat{B} - b\right)\psi\right\rangle\right.
		= \left\langle\psi\left|\left(\hat{A} - a\right)\left(\hat{B} - b\right)\right|\psi\right\rangle \\
		& = \left\langle\psi\left|\hat{A}\hat{B}\right|\psi\right\rangle - b\left\langle\psi\left|\hat{A}\right|\psi\right\rangle 
		- a\left\langle\psi\left|\hat{B}\right|\psi\right\rangle + ab\left\langle\psi\left|\psi\right\rangle\right. \\
		& = \left\langle\hat{A}\hat{B}\right\rangle - ab - ab + ab = \left\langle\hat{A}\hat{B}\right\rangle - ab.
	\end{aligned}\]
	Similarly, $\langle g|f \rangle = \langle\hat{B}\hat{A}\rangle - ab$.\\
	Therefore, 
	\[\langle f|g \rangle-\langle g|f \rangle = \left\langle\hat{A}\hat{B}\right\rangle - \left\langle\hat{B}\hat{A}\right\rangle 
	= \left\langle\left[\hat{A},\hat{B}\right]\right\rangle\]
	Put them all together and we finally get the \emph{Heisenberg uncertainty principle}:
	\[\sigma_{A}\sigma_{B} \geq \left|\frac{1}{2i}\left\langle\left[\hat{A}, \hat{B}\right]\right\rangle\right|.\]
	By letting $\hat{A} = \hat{x}$ and $\hat{B} = \hat{p}$, we obtain the \emph{position-momentum uncertainty principle}:
	\[\Delta x \Delta p \geq \frac{\hbar}{2}.\]
	
	\subsection{Ehrenfest Theorem and Energy-Time Uncertainty Principle}
	Given a time-\emph{dependent} wavefunction $\psi(x, t)$ and an observable $Q$ that \emph{may} depend on time, 
	consider the change of its expectation value with respect to time:
	\[\begin{aligned}
		\frac{\text{d}}{\text{d}t}\langle Q \rangle
		& = \frac{\text{d}}{\text{d}t} \left\langle \psi \left| \hat{Q} \right| \psi \right\rangle \\
		& = \left\langle \frac{\partial\psi}{\partial t} \left| \hat{Q} \right| \psi \right\rangle +
			\left\langle \psi \left| \frac{\partial\hat{Q}}{\partial t} \right| \psi \right\rangle +
			\left\langle \psi \left| \hat{Q} \right| \frac{\partial\psi}{\partial t} \right\rangle
	\end{aligned}\]
	Invoking the time-dependent \emph{Schr\"odinger equation} shown as below:
	\[i\hbar\frac{\partial\psi}{\partial t} = \hat{H}\psi,\]
	where $\hat{H}$ is the \emph{Hamiltonian operator} that represents the whole energy of the state, 
	and note that when taking a constant out of the bra notation, it results in its conjugate.
	\[\begin{aligned}
		\frac{\text{d}}{\text{d}t}\langle Q \rangle
		& = \left(\frac{1}{i\hbar}\right)^*\left\langle \hat{H}\psi \left| \hat{Q} \right. \psi \right\rangle + 
		\left\langle \psi \left| \frac{\partial\hat{Q}}{\partial t} \right| \psi \right\rangle + 
		\frac{1}{i\hbar}\left\langle \psi \left. \hat{Q} \right| \hat{H}\psi \right\rangle\\
		& = -\frac{1}{i\hbar}\left\langle \psi \left| \hat{H}\hat{Q} \right| \psi \right\rangle + 
		\frac{1}{i\hbar}\left\langle \psi \left| \hat{Q}\hat{H} \right| \psi \right\rangle + 
		\left\langle \psi \left| \frac{\partial\hat{Q}}{\partial t} \right| \psi \right\rangle
	\end{aligned}\]
	Cleaning up a bit, and we obtain the \emph{Ehrenfest theorem}:
	\[\frac{\text{d}}{\text{d}t}\langle Q \rangle = \frac{i}{\hbar}\left\langle\left[\hat{H}, \hat{Q}\right]\right\rangle + \left\langle\frac{\partial\hat{Q}}{\partial t}\right\rangle.\]
	It implies that: 1) if $\hat{Q}$ itself doesn't depend explicitly on time, and 2) it commutes with $\hat{H}$, then $Q$ is a \emph{conserved observable}.\\ \\
	Utilizing Heisenberg uncertainty, we set $\hat{A} = \hat{H}$ and $\hat{B} = \hat{Q}$. If $\hat{Q}$ doesn't depend explicitly on time, 
	\[\sigma_{H}\sigma_{Q} \geq \left| \frac{1}{2i}\frac{\hbar}{i}\frac{\text{d}\langle Q \rangle}{\text{d}t} \right| 
	= \frac{\hbar}{2}\left|\frac{\text{d}\langle Q \rangle}{\text{d}t} \right|. \]
	Setting $\Delta E$ equal to $\sigma_{H}$ and $\Delta t$ the time it takes for $\langle Q \rangle$ to change one standard deviation, i.e.
	\[\Delta E = \sigma_{H}\ \ \text{    and    }\ \ \Delta t = \frac{\sigma_{Q}}{\left|\text{d}\langle Q \rangle/\text{d}t\right|},\]
	we obtain the \emph{energy-time uncertainty principle}:
	\[\Delta E \Delta t \geq \frac{\hbar}{2}.\]
	
	\section{To Conclude...}
	The main idea of quantum mechanics covered in two sentences:\\ \\
	The state of a particle is represented by a vector in an inner product space 
	where its properties, i.e. the observables it has, are represented by Hermitian operators 
	of which all possible observable states form a set of orthonormal eigenbasis with the corresponding observable their eigenvalues.
	When observing a quantum state, it will randomly collapse into one of the eigenstates of the observable operators, 
	of which the probability is proportional to the modulus square of the state's component on the eigenstate, 
	and giving the eigenvalue of the eigenstate as a result.